\section{Finite monomial interpolation}

\begin{theorem}[Evaluation recovery of monomial coefficients]
    \label{thm:finite_monomial_interpolation}
    \lean{MvPolynomial.exists_monomial_evaluation_recovery,
        MvPolynomial.monomial_evaluation_recovery_smul}\leanok
    Let $\sigma$ be a set of variables and let $S$ be a finite set of distinct
    monomials. Write $t=|S|$ and $m_\alpha(z)=\prod_i z_i^{\alpha_i}$,
    with only finitely many nonzero exponents for each $\alpha\in S$.
    There exist points $z_0,\ldots,z_{t-1}\in\mathbb C^\sigma$ and complex
    coefficients $c_{\alpha r}$, depending only on $S$, such that
    \begin{align}
        \sum_{r=0}^{t-1}c_{\alpha r}m_\beta(z_r)
         & =\begin{cases}1,&\alpha=\beta,\\0,&\alpha\ne\beta.\end{cases}\notag
    \end{align}
    For every complex vector space $V$ and family $(f_\beta)_{\beta\in S}$ in $V$,
    put $P_f(z)=\sum_{\beta\in S}m_\beta(z)f_\beta$. The same points and
    coefficients satisfy
    \begin{align}
        f_\alpha & =\sum_{r=0}^{t-1}c_{\alpha r}P_f(z_r).\notag
    \end{align}
    Empty $S$ and dependent or zero coefficient vectors are allowed.
    This is the interpolation step in \cite[Lemma~4.2]{Arad2013AreaLaw}.
    The spanning-family interpretation and a separate correction to the
    subsequent word count are recorded in
    \cite{gap:aklv13_interpolation_and_word_count}.
\end{theorem}
\begin{proof}\leanok
    Scalar polynomial evaluation is injective over the infinite field
    $\mathbb C$. The monomial basis therefore gives linear independence of
    the functions $z\mapsto m_\beta(z)$.
    Consequently the evaluation vectors
    $E_z=(m_\beta(z))_{\beta\in S}$ span $\mathbb C^S$.
    Choose $t$ of them spanning this $t$-dimensional space.
    Expressing each coordinate vector as a linear combination of those
    evaluations gives the coefficients $c_{\alpha r}$.
    For vector coefficients, distribute scalar multiplication and interchange
    the two finite sums; the scalar recovery identity leaves precisely $f_\alpha$.
\end{proof}


\begin{theorem}[Rank of a selected coefficient sum]
    \label{thm:monomial_evaluation_selected_rank}
    \lean{Matrix.operatorSchmidtRank_sum_le_of_monomial_evaluations}
    \leanok
    Let $S$ be a finite monomial support, let $f_\alpha$ be complex bipartite
    operator coefficients, and put
    $P_f(z)=\sum_{\alpha\in S}m_\alpha(z)f_\alpha$.
    If $\operatorname{OSR}(P_f(z))\le r$ for every $z\in\mathbb C^\sigma$,
    where $r\ge0$ is an integer, then every $U\subseteq S$ satisfies
    \begin{align*}
        \operatorname{OSR}\!\left(\sum_{\alpha\in U}f_\alpha\right)\le |S|r.
    \end{align*}
    Empty supports, empty subsets, and dependent coefficient operators are allowed.
\end{theorem}
\begin{proof}
    \leanok
    \uses{thm:finite_monomial_interpolation, thm:operator_schmidt_rank_algebra}
    Choose the $t=|S|$ recovery points and coefficients from $S$ before
    choosing the operator family or the selected subset. Summing the recovery
    formulas over $U$ gives
    \begin{align*}
        \sum_{\alpha\in U}f_\alpha
            =\sum_{j=0}^{t-1}\left(\sum_{\alpha\in U}c_{\alpha j}\right)P_f(z_j).
    \end{align*}
    Scalar multiplication does not increase operator Schmidt rank, and the
    rank of a sum is at most the sum of the ranks. The right side therefore
    has rank at most $tr$.
\end{proof}
